% generated by roam.py handbook export -- не редактировать, править заметки
% 20261003154311-сходимость_градиентного_спуска_для_функции_с_липшицевым_градиентом.org (body b801e01ac5ae)

\begin{theorem}[{Сходимость градиентного спуска для функций с липшицевым градиентом}]\label[theorem]{note:C3B64D64-D35B-42AC-B683-6736227CBBBC}
  Пусть $f: \R^n \to \R$, $f \in C_L^{1,1}$~--- \hyperref[note:7DD1C7AB-1A49-4675-9F21-5A16C717F670]{Функции с липшицевым градиентом}, и $f^{*} = \inf_x f(x) > - \infty$.
  Возьмем $\alpha_k = \frac{1}{L}$, а $x_{k+1} = x_k - \alpha_k \nabla f (x_k)$ (см. \hyperref[note:BEEE81AD-13FC-4846-A160-C72E31760BEE]{Градиентный спуск})

  Тогда: Градиентный спуск сходится~--- означает, что $\| \nabla f(x_k)\| \rightarrow 0$, причем
  \[
  \min_{0 \le i \le k} \| \nabla f(x_i)  \| \le \sqrt{\frac{2L (f(x_0) - f^{*})}{k+1}}
  \]
\end{theorem}

\begin{proof}
  По \hyperref[note:FFEA1F1D-B70C-4B4A-9386-8344A2A3BBC4]{оценке убывания за шаг} при $x = x_k$, $y = x_{k+1} = x_k - \alpha \nabla f(x_k)$:
  \[
  f(x_{k+1}) \le f(x_k) - \alpha \left( 1 - \frac{L\alpha}{2} \right) \|\nabla f(x_k)\|^2
  \]
  Подставим $\alpha = \frac{1}{L}$ (Тогда оно максимизирует $\alpha \left(1- \frac{L \alpha}{2}\right)$) \\
  Тогда
  \[
  f(x_{k+1}) \le f(x_k) - \frac{1}{L} \left( 1 - \frac{1}{2}  \right) \|\nabla f(x_k)\|^2 = f(x_k) - \frac{1}{2L} \| \nabla f(x_k)\|^2
  \]
  Давайте перенесем норму налево, а $f(x_{k+1})$ направо
  \[
  \frac{1}{2L} \|\nabla f(x_k)\|^2 \le f(x_k) - f(x_{k+1})
  \]

  Заметим, что $f(x_k) - f(x_{k+1})$ образуют телескопическую сумму:
  \[
  \sum\limits_{k=0}^n [f(x_k) - f(x_{k+1})] = f(x_0) - f(x_{n+1})
  \]
  Также заметим, что $f(x_{n+1}) \ge f^{*}$, он не меньше минимума (естественно); Тогда у нас есть такая оценка
  \[
  \sum\limits_{k=0}^n [f(x_k) - f(x_{k+1})] = f(x_0) - f(x_{n+1}) \le f(x_0) - f^{*}
  \]

  А еще мы можем это оценить снизу: через сумму норм. Тогда
  \[
  \frac{1}{2L} \sum\limits_{k=0}^n \|\nabla f(x_k)\|^2 \le \sum\limits_{k=0}^n [f(x_k) - f(x_{k+1})] \le f(x_0) - f^{*}
  \]

  Ряд $\sum\limits_{k=0}^{\infty} \|\nabla f(x_k)\|^2$ сходится (частичные суммы ограничены числом $2L(f(x_0) - f^{*})$), значит его члены стремятся к нулю: $\|\nabla f(x_k)\| \to 0$. Это первая часть утверждения. Для оценки скорости
  давайте оценим сумму норм градиентов снизу, как $n+1$ минимальных градиентов
  \[
  \frac{n+1}{2L} \min_{0 \le k \le n} \| \nabla f(x_k)\|^2 \le f(x_0) -f^{*}
  \]
  Давайте оставим только минимум слева
  \[
  \min_{0 \le k \le n} \|\nabla f(x_k)\| \le \sqrt{\frac{2L (f(x_0) - f^{*})}{n+1}}
  \]

  Правая часть стремится к нулю как $\frac{c}{\sqrt{n}}$~--- оценка \hyperref[note:A74C35E2-DDF7-4115-99B7-2448CF7C6878]{сублинейного} порядка (быстрее ли сходимость на самом деле, теорема не говорит).

  \begin{remark}
    Оценка $c/\sqrt{n}$ получена для минимума нормы по первым итерациям, а не для $\|\nabla f(x_n)\|$: сама норма градиента стремится к нулю, но не обязана убывать монотонно.
  \end{remark}
\end{proof}
